% concatenated from On__a__Thm__of__Nikol_skii.tex
%------------------------------------------------------------------------------
% Here please write the date of submission of paper or its revisions:
%------------------------------------------------------------------------------
%
\documentclass[12pt, reqno]{amsart}
\usepackage{amsmath, amsthm, amscd, amsfonts, amssymb, graphicx, color}
\usepackage[bookmarksnumbered, colorlinks, plainpages]{hyperref}
\usepackage{float}
\usepackage{lineno}
%\linenumbers
\usepackage{graphicx,epstopdf}
\usepackage{bbm}
\usepackage{mathrsfs}
\usepackage{dsfont}
\usepackage{amssymb}
\usepackage{inputenc}
\usepackage{csquotes}
\usepackage[dvips]{epsfig}
\usepackage{latexsym}
\usepackage{exscale}
%\usepackage[latin1]{inputenc}
%\usepackage{setspace}
%\doublespacing
%\hypersetup{colorlinks=true,linkcolor=red, anchorcolor=green, citecolor=cyan, urlcolor=red, filecolor=magenta, pdftoolbar=true}
%\textheight 22.5truecm \textwidth 14.5truecm
%\setlength{\oddsidemargin}{0.35in}\setlength{\evensidemargin}{0.35in}
%\setlength{\topmargin}{-.5cm}

\newtheorem{theorem}{Theorem}[section]
\newtheorem{lemma}[theorem]{Lemma}
\newtheorem{proposition}[theorem]{Proposition}
\newtheorem{corollary}[theorem]{Corollary}
\theoremstyle{definition}
\newtheorem{definition}[theorem]{Definition}
\newtheorem{example}[theorem]{Example}
\newtheorem{exercise}[theorem]{Exercise}
\newtheorem{conclusion}[theorem]{Conclusion}
\newtheorem{conjecture}[theorem]{Conjecture}
\newtheorem{criterion}[theorem]{Criterion}
\newtheorem{summary}[theorem]{Summary}
\newtheorem{axiom}[theorem]{Axiom}
\newtheorem{problem}[theorem]{Problem}
\theoremstyle{remark}
\newtheorem{remark}[theorem]{Remark}
\numberwithin{equation}{section}

\begin{document}

\setcounter{page}{1}

%\title[Property(E), Reflexivity and Bernstein Lethargy Theorem ]{ Property (E), Reflexivity and Bernstein Lethargy Theorem }
\title{On a theorem of Nikol'skii}

\author[A. AKSOY, Q. Peng]{Asuman G\"{u}ven AKSOY$^*, $ Qidi Peng$^1$}

\address{$^{*}$Department of Mathematics, Claremont McKenna College, 850 Columbia Avenue, Claremont, California  91711, USA.}
\email{\textcolor[rgb]{0.00,0.00,0.84}{aaksoy@cmc.edu}}
\address{$^{1}$Institute of Mathematical Sciences, Claremont Graduate University, 1237 N. Dartmouth Avenue, Claremont, California  91711, USA.}
\email{\textcolor[rgb]{0.00,0.00,0.84}{qidi.peng@cgu.edu}}

%\dedicatory{This paper is dedicated to Professor ABCD}

\subjclass[2010]{Primary 41A25; Secondary 41A50, 46B20.}

\keywords{Best approximation, Bernstein's lethargy theorem, reflexive Banach space.}

%\date{Received: xxxxxx; Revised: yyyyyy; Accepted: zzzzzz.
%\newline \indent $^{*}$Corresponding author}

\begin{abstract}
We present Bernstein lethargy theorem  and  examine the relationship between  Bernstein lethargy theorem and reflexivity.
%%%%%%%%%%%%%%
\iffalse

The equivalence of existence of best approximation  and the reflexivity is well known, moreover equivalence of reflexivity and Bernstein lethargy theorem was proved for a finite sequence of nonnegative numbers with the property 
$e_1\geq e_2\geq \dots\geq e_N$ and $e_{N+1}=0$ for $n\geq N$. In this paper, we consider existence of best approximation together with a sequence of nonnegative numbers  $\{d_n\}$ satisfying

% \label{condition_Borodin}
 $$ d_n >  \sum_{k = n+1}^\infty d_k \quad \mbox{for all}\quad  n \ge n_0 \quad \mbox{ at which}\quad  d_n > 0 $$
  %\end{equation}
  and investigate the relationship with Benstein lethargy theorem and reflexivity.
\fi
%%%%%%%%%%%%
\end{abstract} \maketitle

\section{Introduction}
%%%%%%%%%%%%
%%%%%%%%%%%%%%%%%
%%%%%%%%%%%%%%%%%%%%
%%%%%%%%%%%%%%%%%%%%%
\subsection{Bernstein Lethargy Theorem}
 One of the notable theorems used in the constructive theory of functions, called the Bernstein's Lethargy Theorem (BLT) \cite{Bernstein}. For a subspace $Y$ of a normed linear space $(X,\|\cdot\|)$,  define the distance from an element $x \in X$ to $Y$ by
$$\rho(x, Y):= \inf \{ \|x-y\|: \,\,y\in Y\}.$$
The Weierstrass Approximation Theorem states that polynomials are dense in $C[0,1]$, thus it is known that for any $f\in C[0,1]$, $$\displaystyle \lim_{n \to \infty} \rho(f, \Pi_n) = 0,$$
where for each $n\ge1$,  $\Pi_n$ is the vector space of all real polynomials of degree at most equal to $n$.
However, the Weierstrass Approximation Theorem gives no information about the speed of convergence for $\rho(f, \Pi_n)$. In $1938$, S.N. Bernstein \cite{Bernstein} proved that if $\{d_n \}_{n\ge1}$ is a non-increasing  sequence  of positive numbers with $\lim\limits_{n\to\infty}d_n=0$, then there exists a function $f\in C[0,1]$ such that
$$\rho(f, \Pi_n)= d_n,~\mbox{for all $n \ge1$}.$$
This remarkable result is called the Bernstein's Lethargy Theorem (BLT) and is used in the constructive theory of functions \cite{Sin}, and  in the theory of quasi-analytic functions in several complex variables \cite{Ple2}.

More generally, let $\{Y_n\}$ be a system of strictly nested subspaces of the Banach space $X$. The sequence of errors of the best approximation from $x\in X$ to $Y_n$, denoted by $\{\rho(x,Y_n)\}$, may converge to zero at an arbitrarily slow rate  or an arbitrarily fast rate. For example, Shapiro in \cite{Sha}, replacing $C[0,1]$ with an arbitrary Banach space $(X, \|\cdot\|)$ and $\{\Pi_n\}$, with a sequence $\{Y_n\}$ where $Y_1 \subset Y_2 \subset \cdots$ are strictly embedded closed subspaces of $X$, showed that in this setting, $\rho(x, Y_n)$ can decay arbitrarily slowly. More precisely, for each null sequence (i.e. a sequence converging to $0$) $\{d_n\}$ of non-negative numbers, there exists $x\in X$ such that
\begin{equation*}
\label{Shap}
 \rho(x, Y_n) \neq \mathcal O (d_n),~\mbox{as $n\to\infty$}.
 \end{equation*}
 This result is further improved  by Tyuriemskih \cite{Tyu}, who showed that for any expanding sequence $\{Y_n\}$ of subspaces of $X$ and for any null sequence $\{d_n\}$ of positive numbers, there is an element $x\in X$  such that
 \begin{equation*}
 \label{tyu}
 \lim_{n \rightarrow \infty} \rho(x, Y_n) =0,~\mbox{and}~\rho(x, Y_n) \geq d_n~\mbox{for all $n\ge1$}.
   \end{equation*}
   However, it is also possible that the errors of the best approximation $\{\rho(x,Y_n)\}$ may converge to zero arbitrarily fast. For results of this type see   \cite{Al-To}.

    We refer the reader to \cite{Ak-Al,Ak-Le,Al-To, Borodin1,Lew1} for other versions of the Bernstein's Lethargy Theorem  and  to \cite{Ak-Le2,Alb,Mich,Ple1,Vas} for the Bernstein's Lethargy Theorem for Fr\'{e}chet  spaces. Consult \cite{De-Hu} for an application of Tyuriemskih's Theorem to convergence of sequence of bounded linear operators and  \cite{Al-Oi} for a generalization of Shapiro's Theorem.
    %%%%%%%%%%%%
    We first  state a  Bernstein Lethargy result concerning \textit{finite number} of subspaces, for the proof  of the following lemma we refer the reader to Timan's book  \cite{Tim}.
\begin{lemma}
\label{lemma1}
Let $(X,\|\cdot\|)$ be a normed linear space, $Y_1\subset Y_2\subset\ldots\subset Y_n\subset X$ be a finite system of strictly nested subspaces, $d_1>d_2>\ldots>d_n\ge0$ and $z \in X\backslash Y_n$. Then, there is an element $x\in X$ for which $\rho(x,Y_k)=d_k$ $(k=1,\ldots,n)$,  $\|x\|\le d_1+1$, and $x-\lambda z\in Y_n$ for some $\lambda>0$. If all $Y_k$ are existence subspaces (in particular, if all $Y_k$ are finite-dimensional or $X$ is reflexive),
then the above assertion holds for arbitrary numbers $d_1 \ge d_2 \ge \ldots\ge d_n \ge 0$.
\end{lemma}

    
    %%%%%%%%%%%

 The general question is that, given a Banach space $X$ and a sequence of strictly increasing subspaces $\{Y_n\}$ of $X$  with a dense union and a decreasing null sequence $ \{d_n\}$, does there exist an $x\in X$ such that $\rho(x, Y_n) =d_n$? If such an element exists for all such subspaces $Y_n$ and numbers $d_n$ , then we say $X$ has the $(B)$-property. Besides Hilbert spaces \cite{Tyu2},  there are no known examples of infinite-dimensional  spaces $X$ having the $(B)$-property.
  Although for a long time no sequence $\{d_n\}$ of this type was known for which such an element $x$ exists for \textit{all} possible Banach spaces $X$, Borodin in \cite{Borodin1} showed that the problem has a solution whenever the sequence $\{d_n\}$ satisfies some fast convergence condition, or the subspaces $\{Y_n\}$ satisfy some  conditions. The first result in \cite{Borodin1}, obtained from the sequence condition (\ref{condition_Borodin}) is stated below:
\begin{theorem}{\cite[p. 624]{Borodin1}}
  \label{thm:Borodin1}
  Let $X$ be an arbitrary infinite-dimensional Banach space, $Y_1 \subset Y_2
  \subset \cdots$ be an arbitrary countable system of strictly nested subspaces
  in $X$, and fix a numerical sequence $\{d_n\}_{n\ge1}$ for which there exists a natural number $n_0\ge1$ such that
  \begin{equation}
  \label{condition_Borodin}
  d_n >
  \sum_{k = n+1}^\infty d_k~\mbox{for all $n \ge n_0$ at which $d_n > 0$.}
  \end{equation}
  Then there is an element $x \in X$ such that
  \begin{equation}
  \label{rhod}
  \rho(x,Y_n) = d_n, ~\mbox{for all $n\ge1$}.
  \end{equation}
\end{theorem}
The condition (\ref{condition_Borodin}) on the sequence $\{d_n\}$ is the key to the derivation of (\ref{rhod}) in Theorem \ref{thm:Borodin1}.  Note that, the condition (\ref{condition_Borodin}) is satisfied when $d_n =(2+\epsilon)^{-n} $ for $\epsilon > 0$ arbitrarily small, however it is not satisfied when $d_n= 2^{-n}$. 
The second result in \cite{Borodin1}, based on the subspace condition (\ref{Borodin2}) is given below:
\begin{theorem}{\cite[p. 626]{Borodin1}}
\label{thm:Borodin2}
	Let $d_0\ge d_1 \geq d_2 \geq \dots > 0$ be a non-increasing sequence converging to $0$ and $Y_1 \subset Y_2 \subset \dots \subset X$ be a system of strictly nested subspaces of  an infinite-dimensional Banach space $X$ for which there exists a series of non-zero elements $\{q_n\}$ such that $q_n \in Y_{n + 1} \setminus Y_n$, and
\begin{equation}
	\label{Borodin2}
\|q\| \leq \displaystyle \frac{d_{k-1}}{d_{k}}\rho(q,Y_k)
\end{equation}
holds for all $k \in \mathbb N$ and any non-zero element $q$ in the linear span $\langle q_k, q_{k + 1},\dots \rangle$. Then there is some element $x$ in the closed linear span $\overline{\langle q_1, q_2,\dots \rangle}$ satisfying $$\rho(x,Y_n) = d_n ~~for ~all ~n \ge1.$$
\end{theorem}
     The latter result is improved by Aksoy et al. in  \cite{Aksoy}. Note that Borodin's proof works only under the assumption that $\overline{Y}_n$ is strictly included  in $Y_{n+1}$ . Necessity of this assumption on subspaces is illustrated by the following example.
    %%%%%%%%%%%%%
  \begin{example}\label{EXAM}
Let $X = L^{\infty}[0,1]$ and consider $C[0,1] \subset L^{\infty} [0,1]$ and define the subspaces  of $X$ as follows:

\begin{enumerate}
\item $Y_1 =  \mbox{space of all polynomials}$;

\item $Y_{n+1} = $span$[Y_n \cup \{f_{n}\}]$
where $f_{n} \in C[0,1] \setminus Y_n$, for $n\ge1$.
\end{enumerate}
Observe that by the Weierstrass Theorem we have $\overline {Y}_n = C[0,1] $ for all $n \ge1$.
Take any  $ f \in L^{\infty}[0,1]$ and consider the following cases:
\begin{enumerate}
\renewcommand{\labelenumi}{\alph{enumi})}
\item If  $ f \in C[0,1]$ , then
$$
\rho (f,Y_n) = \rho(f,C[0,1]) = 0~\mbox{for all $n\ge1$}.
$$
\item If $ f \in L^{\infty}[0,1] \setminus C[0,1] $, then
$$
\rho(f, Y_n) = \rho(f, C[0,1]) = d > 0~ \mbox{ (independent of $n$)}.
$$
Note that in the above, we have used the fact that $\rho(f,Y_n) = \rho(f,\overline{Y}_n)$.
Hence in this case Bernstein lethargy theorem does not hold.
\end{enumerate}
\end{example}
  
    
    %%%%%%%%%%%%%
      It is natural to ask the relationship between reflexivity and Bernstein Lethargy Theorem. After all, reflexive spaces have better approximation properties, such as the existence of bounded projections or the weak compactness of closed balls. Moreover, certain versions of Bernstein Lethargy Theorem require the space itself or subspaces to be reflexive. We examine this in more detail in the next section.

%%%%%%%%%%%%%%
%%%%%%%%%%%%%
%%%%%%%%%%%%
 %%%%%%%%%%%%%%
\subsection{Existence Property and Reflexivity}

Let $X$ be a real or complex Banach space and $Y$ be a closed linear subspace of $X$. An element $y_0\in Y$ is called the \textit{the best approximation or the least deviation} from $x\in X$ if $$||x-y_0||=\inf_{y\in Y} ||x-y||=\rho(x, Y).$$ 
We say a Banach space $X$ has the existence property or \textit{Property (E)} for short, if for every choice of closed subspace $Y$ and the vector $x\in X$ there is at least one best approximation $y_0$  for $x$ in $Y$. 
\begin{theorem}\cite{Hir}\label{Hir}
Let $X$ be a Banach space, then $X$ has property $(E)$ if and only if $X$ is reflexive.
\end{theorem}
If $X$ is reflexive then any closed linear subspace $Y$ of $X$ is reflexive and from the fact that  in a reflexive space, the  closed unit ball is weakly compact, shows that $X$ has the property $(E)$.  The other implication comes from the result of R.C. James in  \cite{RCJ}, where he proves that if every linear functional assumes its maximum on absolute value on the unit sphere of any closed subspace of $X$, then $X$ is reflexive.
Another  equivalent way of looking at reflexive spaces are through the concept of proximal subspaces.  We call  $Y_n$ \textit{proximinal} in $X$ if for any $x\in X$ there exists $v\in Y_n$ such that
 $$\rho(x, Y_n) =\|x-v\|.$$ By letting $y=x-v$ we obtain that if    $Y_n$ is  proximinal in  $Y_{n+1}$, then given $x\in Y_{n+1}$,  there is $v\in Y_n$ such that $\rho(x,Y_n)=\|x-v\|$ , this condition reduces to for each $n\ge1$, there exists $y\in Y_{n+1}\backslash Y_n$ such that
\begin{equation}
\label{Yn'}
\rho(y, Y_n)=\|y\|
\end{equation}

 Thanks to the characterization of reflexivity in terms of the proximinality of its closed subspaces, the condition (\ref{Yn'})  holds for reflexive Banach spaces (see \cite[Theorem $2.14$, p. $19$]{Singer2}). \\
 
 Next we list some properties of the distance function $\rho(x, Y_n)$:\\[.02in]
Given a Banach space $X$ and its subspaces $Y_1\subset Y_2 \subset \dots \subset Y_n\subset \cdots$, it is clear that $$ \rho(x, Y_1) \geq\rho(x,Y_2)\geq \cdots, ~\mbox{for any $x\in X$}$$ and  thus $\{\rho(x,Y_n)\}_{n\ge1}$ forms a non-increasing sequence of errors of best approximation from $x$ to $Y_n$, $n\ge1$.  Furthermore,  it is easy to show that
\begin{eqnarray*}
&&\rho(\lambda x, Y_n)= |\lambda| \rho(x, Y_n),~\mbox{for any}~ x\in X~\mbox{and}~\lambda\in\mathbb R;\\
&& \rho(x+v, Y_n)= \rho(x, \overline Y_n),~\mbox{for any}~x\in X~ \mbox{and}~v\in \overline Y_n;\\
&&\rho(x_1+x_2, Y_n)\leq \rho(x_1, Y_n)+\rho(x_2, Y_n)
\end{eqnarray*}
and consequently   $$\rho(x_1+x_2, Y_n)\geq |\rho(x_1, Y_n)-\rho(x_2, Y_n)|,~ \mbox{ for any}~  x_1,x_2\in X.$$
 Note that we also have: $$ |\rho(x_1, Y_n)-\rho(x_2, Y_n)| \leq \|x_1-x_2\|,~ \mbox{ for any}~  x_1,x_2 \in X,$$ which  implies that  the mapping $X \longrightarrow \mathbb{R}^+$  defined by $ x\longmapsto \rho(x, Y_n)$ is  continuous and thus properties of continuous mappings such as the intermediate value theorem can be used. Moreover 
 $ \rho(x+ty, Y_n)$ is a continuous function of $t$ and
 if $y \notin Y_n$ then $ \rho(x+ty, Y_n) \to \infty$ as $|t| \to \infty$.

\subsection{Bernstein Lethargy Theorem and Reflexivity}


\begin{definition}\cite{Nik}
A Banach space $X$  has the property $(B_f)$ if for every countable sequence of closed  distinct subspaces 
$$\{0\} = Y_0\subset Y_1\subset \dots\subset Y_n\subset \dots\subset X$$ and any finite sequence of nonnegative numbers
$$ e_0\geq e_1\geq \dots\geq e_N$$ there exists $x\in X$ for which $\rho(x,Y_n)=e_n$ for $n=1,2,\dots, N$ and $\rho(x,Y_n)=0$ when $n=N+1,\dots$. 

\end{definition}
\begin{theorem}\cite{Nik}
\label{B_f}
A Banach space  $X$ has the property $(B_f)$ if and only if $X$ is reflexive. 
\end{theorem}
Since the paper   \cite{Nik}  of Nikol'skii is  not available  in translation and appeared in a hard to reach regional  journal, in the following we give the proof by clarifying the steps.
\begin{proof}
Suppose $X$ is not reflexive,  there exists $Y$, non-reflexive subspace of $X$, such that its complementary $Y^c:={\rm span}(X\setminus Y)$ is infinite dimensional, and we use this to construct a sequence of closed subspaces $Y_0=\{0\}\subset \cdots \subset Y_n \subset Y_{n+1}\subset \cdots\subset X$ with certain properties. In particular, there is $y_0\in X$ such that $\rho(y_0,Y_1)=\rho_0>0$ but there are no best approximations to $y_0$ with the elements of $Y_1$. In particular, $\|y_0\|_X>\rho_0$ since otherwise $0$ would be a best approximation to $y_0$ from $Y_1$.  Moreover, $Y_2=\overline{{\rm span}\{y_0,Y_1\}}^X$. You claim that there is no $x\in X$ such that 
\[
\|x\|=\rho(x,Y_0)=\rho (x,Y_1)=\rho_0>0=\rho(x,Y_2)=\rho(x,Y_3)=\cdots 
\]
Clearly, $0=\rho(x,Y_2)$ means that $x\in Y_2=\overline{{\rm span}\{y_0,Y_1\}}^X$. If $x=\lambda y_0+y_1$ with $y_1\in Y_1$ and $\lambda\in \mathbb{R}$ , then $\rho(x,Y_1)=\rho(\lambda y_0,Y_1)=|\lambda|\rho_0=\rho_0$ implies $|\lambda|=1$, so $x=\pm y_0+y_1$. Then $0=\|x\|=\rho(x,Y_0)$ would imply $$\|\pm y_0+y_1\|=\|y_0\pm y_1\|=\rho_0=\rho(y_0,Y_1)$$ and $y_0$ would admit a best approximation from $Y_1$, which is impossible. 

Now, it may happen that $x=\lim_{n\to\infty}(\lambda_ny_0+y_n)$ for certain sequences $(\lambda_n)\subseteq \mathbb{R}$ and $(y_n)\subseteq Y_1$, but $x$ does not belong to ${\rm span}\{y_0,Y_1\}$.  In such a case, we would get
$$\rho(x,Y_1)=\lim_{n\to\infty}\rho
(\lambda_n y_0,Y_1)=\lim_{n\to\infty}|\lambda_n|\rho_0=\rho_0,$$ which just means $\lim_{n\to\infty}|\lambda_n|=1$. 

%We need to get a contradiction from these facts and the fact that 
%\begin{equation} \label{uno}
%    \rho_0=E(x,Y_0)=\|x\|=\lim_{n\to\infty} \|\lambda_ny_0+y_n\|_X. 
%\end{equation}

Now, $\lim_{n\to\infty}|\lambda_n|=1$ implies that we can extract a subsequence $(\lambda_{n_k})$ from $(\lambda_n)$ that converges either to $1$ or to $-1$ (this is so because we only deal with real numbers). Assume that $\lim_{k\to\infty}\lambda_{n_k}=1$. Then $x=\lim_{k\to\infty}(\lambda_{n_k}y_0+y_{n_k})= y_0+\lim_{k\to\infty}y_{n_k}$ means that $\lim_{k\to\infty}y_{n_k}=x-y_0$ exists, so it belongs to $Y_1$ (which is closed). A contradiction since we assumed that $x$ does not belong to ${\rm span}\{y_0,Y_1\}$! An analogous argument applies if $\lim_{k\to\infty}\lambda_{n_k}=-1$. 

%%%%%%%%%%%%%%%%%%%
\iffalse
then there exists a non-reflexive subspace $Y$  such that $X\setminus Y$ is infinite dimensional.  Based on the Theorem \ref{Hir} there is an element $y_0\in Y$ and a closed subspace $Y^* \subset Y$ such that for $y_0$ there does not exist in $Y^*$ a least deviation.  Suppose $\rho(y_0, Y^*)= \rho_0> 0$. Next set $Y_1=Y^*$ and let $Y_2 $ be the closed linear span of $y_0$ and $Y^*$. In this manner we can define a countable sequence $\{Y_n\}_{n=3}^{\infty}$ where all $Y_n$'s are distinct closed subspaces of $X$. This is possible since $X\setminus Y$ is infinite dimensional.  Now, define numbers $e_n$ with
$$ e_0=e_1=\rho_0, \quad \quad e_n=0\quad\mbox{for}\quad n\geq 2.$$  $y_0$ is the unique element up to constant $\epsilon$, $|\epsilon|=1$ that satisfies the following conditions:
$$ y_0\in Y_2 \quad \mbox{and} \quad  \rho(y_0, Y_1)=e_1=\rho_0$$
The equation $\rho( y_0, Y_0)=  e_0= \rho_0$ can not hold. Since in this case there is an element  in $Y_1$  that will be the least deviation from $y_0$. Thus there does not exist $Y_1=Y^*$ with the $\rho(y_0, Y_1)=e_1=\rho_0$.  
\fi
%%%%%%%%%%%%%%%%%%%%%%%%%%
i.e., if $X$ is not reflexive  then $X$ does not have the property $(B_f)$.


To prove the other implication, suppose $X$ is reflexive, take any element $y_{N+1} \in Y_{N+1} \setminus Y_N$, then find $t=c_{N+1}$ such that $$ \rho( c_{N+1} y_{N+1}, Y_N)= e_N.$$ Since $X$ is reflexive, there exists an element $y\in Y_N$ which is the least deviation from $ c_{N+1} y_{N+1}$.

Hence we have $$ \rho( c_{N+1} y_{N+1}+ y , Y_0)= e_N.$$

Let $y\in Y_m\setminus Y_{m-1}$, consider the following possibilities:
%%%%%%%%%%%%%%%%
\begin{enumerate}
\renewcommand{\labelenumi}{\alph{enumi})}

%\begin{description}
\item$m = N$ and $e_N \le e_{N-1}$\\


In this case, take $y_N\in Y_N\backslash Y_{N-1}$ to be the least deviation from $c_{N+1}y_{N+1}$ in $Y_N$,
$$
 \rho(c_{N+1}y_{N+1}-y_N, Y_{N-1} )\le \|c_{N+1}y_{N+1}-y_N\|=e_N\le e_{N-1}.
$$
Also since $\rho(y_N, Y_{N-1} )>0$, by the triangle inequality of $\rho(\cdot,Y)$, 
$$
 \rho(c_{N+1}y_{N+1}+\lambda y_N, Y_{N-1} )\ge  \lambda \rho(y_N, Y_{N-1} )- \rho(c_{N+1}y_{N+1}, Y_{N-1} )\xrightarrow[\lambda\to+\infty]{}+\infty.
$$
Therefore by the intermediate value theorem, there is $c_{N}\in\mathbb R$ such that
$$
\rho(c_{N+1}y_{N+1}+c_N y_N, Y_{N-1} )=e_{N-1}.
$$
\item $m < N$ and $e_N = e_{N-1}$\\

In this case, taking $c_N=1$, we directly get
$$
\rho(c_{N+1}y_{N+1} + c_Ny_N, Y_N ) = \rho(c_{N+1}y_{N+1} + y_N, Y_N )=e_N=e_{N-1}.
$$
\item $m < N$ and $e_N < e_{N-1}$\\

In this case, for any $y\in Y_N\backslash Y_{N-1}$, 
$$
\|c_{N+1}y_{N+1}-y\|>\rho(c_{N+1}y_{N+1},Y_N)=e_N.
$$
Let us denote $\|c_{N+1}y_{N+1}-y\|=e_N+\delta$ with $\delta>0$. Then by the triangle inequality,
$$
\|c_{N+1}y_{N+1}-cy\|\ge \|c_{N+1}y_{N+1}-y\|-|c-1|\|y\|=e_n+\delta-|c-1|\|y\|.
$$
For any $\epsilon\in(0,e_{N-1}-e_N)$, take $c\in\mathbb R$ such that $\delta-|c-1|\|y\|=\epsilon$ and $y_N=cy$.  
Then $y_N\in Y_N\backslash Y_{N-1}$ and
$$
\rho(c_{N+1}y_{N+1}-y_N,Y_{N-1})\le \|c_{N+1}y_{N+1}-y_N\|=e_N+\epsilon<e_{N-1}.
$$
Also since $\rho(y_N, Y_{N-1} )>0$, by the triangle inequality of $\rho(\cdot,Y)$, 
$$
 \rho(c_{N+1}y_{N+1}+\lambda y_N, Y_{N-1} )\ge  \lambda \rho(y_N, Y_{N-1} )- \rho(c_{N+1}y_{N+1}, Y_{N-1} )\xrightarrow[\lambda\to+\infty]{}+\infty.
$$
Therefore by the intermediate value theorem, there is $c_{N}\in\mathbb R$ such that
$$
\rho(c_{N+1}y_{N+1}+c_N y_N, Y_{N-1} )=e_{N-1}.
$$
\end{enumerate}
%%%%%%%%%%%%%%%%%
%%%%%%%%%%%%%%

In all these three cases, using  the reflexivity of the subspaces  and the properties  $\rho\,\,(\,\, .\,\, , Y)$,  in particular continuity  it can be shown that $$\rho(c_{N+1} y_{N+1}+c_Ny_N, Y_{N-1} )= e_{N-1}.$$ Continuing in this manner we can construct an element $$x =c_{N+1} y_{N+1}+\dots+c_1y_1$$ which has the desired properties.
\end{proof}
 Note that Nikol'skii's result is  restricted to a special sequence of the form
 $$d_0\geq d_1\geq \dots  \geq d_N \geq d_{N+1} =0 =d_{N+m}$$ for all $m\geq 1$.\\
 
  The general question is that, given a Banach space $X$ and a sequence of strictly increasing subspaces $\{Y_n\}$ of $X$  with a dense union and a decreasing null sequence $ \{d_n\}$, does there exist an $x\in X$ such that $\rho(x, Y_n) =d_n$? If such an element exists for all such subspaces $Y_n$ and numbers $d_n$ , then we say $X$ has the $(B)$-property. Besides Hilbert spaces \cite{Tyu2}, ( see \cite{AL} for an alternative proof)  there are no known examples of infinite-dimensional  spaces $X$ having the $(B)$-property. Nevertheless one can show the following relationship between the property (B) and the reflexivity.

\begin{theorem}
If $X$ satisfies the property (B) then X is reflexive.
\end{theorem}
\begin{proof}
Suppose   $d_1 = d_2=1$, for $n \geq 3$  the sequence $\{d_n\}$  is non-increasing and $\{d_n\} \to 0$.
Consider a pair of subspaces $Y_1\subset Y_2 \subset X$ such that
$$
Y_1= \{0\};\quad \quad Y_2= \ker (f)= \{ y\in X: \,\, f(y)=0\},
$$
where $f$ is a nonzero linear functional. Suppose the Bernstein's Lethargy Theorem holds.  In other words, there is an element $x\in X$ with prescribed best approximation errors. Since in this case

 $$ \rho(x, Y_1)= \rho(x, Y_2)= 1$$ and

$$\rho(x, Y_1)= \|x\|, \quad \rho(x, Y_2)= \displaystyle\frac{f(x)}{\|f\|}, $$  one has the equalities
$$ 1= \|x\|= \displaystyle\frac{f(x)}{\|f\|}.$$
 From the last equality we see that the functional $f$ assumes its supremum at $x$.   This implies that  $X$ is reflexive (see James' Theorem \cite{RCJ}).
\end{proof}

\begin{remark}
The next  step is to consider the existence of best approximation together with a sequence of nonnegative numbers $d_n$ satisfying $\displaystyle \lim_{n\to \infty} d_n=0$ and investigate the relationship between Bernstein Lethargy Theorem and reflexivity. 
\end{remark}
 %%%%%%%%%%%%%%%
 \iffalse
 
 
 %%%%%%%%%%%%%%%%%%%%%
 %%%%%%%%%%%%%%%%%%%%%%%
  %%%%%%%%%%%%%%%%%
  \iffalse
    Although for a long time no sequence $\{d_n\}$ of this type was known for which such an element $x$ exists for \textit{all} possible Banach spaces $X$, Borodin in \cite{Borodin1} showed that the problem has a solution whenever the sequence $\{d_n\}$ satisfies some fast convergence condition, or the subspaces $\{Y_n\}$ satisfy some  conditions. The first result in \cite{Borodin1}, obtained from the sequence condition (\ref{condition_Borodin}) is stated below:
\begin{theorem}{\cite[p. 624]{Borodin1}}
  \label{thm:Borodin1}
  Let $X$ be an arbitrary infinite-dimensional Banach space, $Y_1 \subset Y_2
  \subset \cdots$ be an arbitrary countable system of strictly nested subspaces
  in $X$, and fix a numerical sequence $\{d_n\}_{n\ge1}$ for which there exists a natural number $n_0\ge1$ such that
  \begin{equation}
  \label{condition_Borodin}
  d_n >
  \sum_{k = n+1}^\infty d_k~\mbox{for all $n \ge n_0$ at which $d_n > 0$.}
  \end{equation}
  Then there is an element $x \in X$ such that
  \begin{equation}
  \label{rhod}
  \rho(x,Y_n) = d_n, ~\mbox{for all $n\ge1$}.
  \end{equation}
\end{theorem}
The condition (\ref{condition_Borodin}) on the sequence $\{d_n\}$ is the key to the derivation of (\ref{rhod}) in Theorem \ref{thm:Borodin1}.  Note that, the condition (\ref{condition_Borodin}) is satisfied when $d_n =(2+\epsilon)^{-n} $ for $\epsilon > 0$ arbitrarily small, however it is not satisfied when $d_n= 2^{-n}$. 
%In this framework we provide an improvement on  Theorem \ref{thm:Borodin1}.

Next, we state a  Bernstein Lethargy result concerning \textit{finite number} of subspaces, for the proof  of the following lemma we refer the reader to Timan's book  \cite{Tim}.
\begin{lemma}
\label{lemma1}
Let $(X,\|\cdot\|)$ be a normed linear space, $Y_1\subset Y_2\subset\ldots\subset Y_n\subset X$ be a finite system of strictly nested subspaces, $d_1>d_2>\ldots>d_n\ge0$ and $z \in X\backslash Y_n$. Then, there is an element $x\in X$ for which $\rho(x,Y_k)=d_k$ $(k=1,\ldots,n)$,  $\|x\|\le d_1+1$, and $x-\lambda z\in Y_n$ for some $\lambda>0$. If all $Y_k$ are existence subspaces (in particular, if all $Y_k$ are finite-dimensional or $X$ is reflexive),
then the above assertion holds for arbitrary numbers $d_1 \ge d_2 \ge \ldots\ge d_n \ge 0$.
\end{lemma}
\fi

%%%%%%%%%%%%%%%%%%
%%%%%%%%%%%%%%%%%%%
%%%%%%%%%%%%%%
%In this framework we provide an improvement on  Theorem \ref{thm:Borodin1}.
%%%%%%%%%%%%%

 We say a Banach space  $X$ has the (BLT)-property if it satisfies the main Theorem \ref{Borodin} below (see Definition \ref{def:BLT}).  This theorem yields an equivalence between Banach spaces  which satisfy  the (BLT)-property and reflexivity.  It is known that if a given Banach space satisfies the (BLT)-property, then it is reflexive. But it was an open problem  for a long time that if $X$ is a reflexive Banach space  whether or not  $X$ satisfies the (BLT)-property (see p. 152 of \cite{Sin}).  The corollary (see Corollary \ref{corollary:main}) to our main Theorem \ref{Borodin} answers this question in the affirmative. Thus, our main result gives a new characterization of reflexivity.

 %%%%%%%%%%%%%%%%%%%%
%%%%%%%%%
%%%%%%%%%%%%%%%%%

%%%%%%%%%%%%%%%
  An element $x\in X$ satisfying $\rho(x, Y_n)=d_n$, $n\ge1$ may exist if the sequence $\{d_n\}$ decreases strictly to zero. Under the above assumption on $\{d_n\}$, Borodin  in \cite{Borodin1} uses Lemma \ref{lemma1}  to establish the existence of such an element $x$ in the case of rapidly decreasing sequences (see Theorem  \ref{thm:Borodin1}). We will assume that the subspaces $\{Y_n\}$ satisfy $\overline Y_n\subset Y_{n+1}$ for $n\ge1$  for the rest of the paper.
  %%%%%%%%%%%%%%%
 %%%%%%%%%%%%%%%%%%%%%%%%%%%%%%%%%%%%%
\section{Main Result}

%%%%%%%%%%%%%%%%
\iffalse
Our first main result gives a positive answer to Question $1$, by showing that Theorem \ref{thm:Borodin} can be extended by weakening the strict inequality in (\ref{condition_Borodin}) to a non-strict one:
% (see (\ref{condition}) below).
\begin{equation}
 \label{condition}
 d_n\ge\sum_{k=n+1}^\infty d_k,~\mbox{for every} \,\, n \ge n_0.
 \end{equation}
%\begin{theorem}
%\label{Borodin}
%Let $X$ be an arbitrary infinite-dimensional Banach space, $\{Y_n\}_{n\ge1}$ be an arbitrary system of strictly embedded subspaces with the property $\overline Y_n \subset Y_{n+1}$ for all $n\ge1$, and let the non-negative numbers $\{d_n\}_{n\ge1}$ satisfy the following property: there is an integer $ n_0 \ge 1 $ such that
 %\begin{equation}
 %\label{condition}
 %d_n\ge\sum_{k=n+1}^\infty d_k,~\mbox{for every} \,\, n \ge n_0.
% \end{equation}
% Then there exists an element $x\in X$ such that $\rho(x,Y_n)=d_n$ for all $n\ge1$.
%\end{theorem}
 Clearly the condition  (\ref{condition}) is weaker than (\ref{condition_Borodin}),  but unlike the condition (\ref{condition_Borodin}), (\ref{condition}) is satisfied by the sequences $\{d_n\}_{n\ge1}$ verifying $d_n=\sum\limits_{k=n+1}^\infty d_k$ for all $n\ge n_0$. For a typical example of such sequence one can take $\{d_n\} =\{2^{-n}\}$. As a consequence, the proof of Theorem \ref{Borodin} requires a finer construction of the element $x$ than that of Theorem \ref{thm:Borodin}. Before proving Theorem \ref{Borodin},
 \fi
 %%%%%%%%%%%%%%%%%

Before proceeding with our main result  we state the following two technical lemmas that we need for the  proof of  the construction of an element $x\in X$ in our main theorem. We include the proof of both of these lemmas in the Appendix.

The first Lemma \ref{lemma1'} 
extends the classical Hahn-Banach Theorem. Unlike the classical Hahn-Banach Theorem (see e.g. \cite[p. 63]{Dunford}), where the linear functional $f$ is evaluated on one element $x$, in  Lemma \ref{lemma1'} the values of $f$ are evaluated on two elements $x_1,x_2$. The second  Lemma \ref{lemma 2}  reveals that, small changes of the values $\{d_n\}$ can lead to small position changes of the corresponding element $x\in X$.


\begin{lemma}
\label{lemma1'}
Let $(X,\|\cdot\|)$ be a normed linear space and let $Q$ be a subspace of $X$ with $\overline Q\subset X$. Pick two elements $x_1,x_2\in X\backslash \overline Q$ such that $x_2\notin \mbox{span} [\{x_1\}\cup Q]$. Then, there exists a non-zero linear functional $f:~X\longmapsto \mathbb R$ such that
\begin{equation}
\label{linearf1}
f(q)=0,~\mbox{for all $q\in Q$},~\|f\|=\frac{1}{\rho(x_1,Q)},~f(x_1)=1
\end{equation}
and
\begin{equation}
\label{linearf2}
f(x_2)=\lim_{a\to+\infty}\left(a-\frac{\rho(x_2-a x_1,Q)}{\rho(x_1,Q)}\right).
\end{equation}
\end{lemma}
%%%%%%%%%%%%%%%%%%%%%


  The following Lemma \ref{lemma2} is another technical statement, it reveals that, small changes of the values $\{d_n\}$ can lead to small position changes of the corresponding element $x\in X$.
 \begin{lemma}
 \label{lemma2}
 Let $Q_1,Q_2$ be two subspaces of an arbitrary normed linear space $(Q_3,\|\cdot\|)$, such that $\overline Q_k\subset Q_{k+1}$ for $k=1,2$. Also assume that for $k=1,2$ there is $y_k\in Q_{k+1}\backslash  Q_k$ for which $\rho(y_k,Q_k)=\|y_k\|$. Let $\{u_m\}_{m\ge1}$ and $\{v_m\}_{m\ge1}$ be two sequences of non-negative numbers, with $u_m\ge v_m$ for all $m\ge1$. Then there exists a sequence of elements  $\{q_m\}_{m\ge1}\subset Q_{3}$, given by
$$
q_{m}=v_m(z-f(z)w)+\mu_mf(z)w,~\mbox{for $m\ge1$},
$$
such that
\begin{enumerate}
\renewcommand{\labelenumi}{\alph{enumi})}
\item [(i)]$z$ is an element in $Q_3\backslash \overline Q_2$ satisfying $\rho(z,Q_1)=\rho(z,Q_2)=1$; $w$ is an element in $\overline Q_3\backslash\{0\}$ such that $\|z-w\|=1$.
    \item[(ii)] $f:~Q_3\rightarrow \mathbb R$ is a non-zero linear functional such that 
    $$
    f(z)=\lim_{a\to+\infty}\left(a-\frac{\rho(z-aw,Q_1)}{\rho(w,Q_1)}\right).
    $$
    \item[(iii)] For each integer $m\ge1$, the variable $\mu_m$ in the expression of $q_m$ is some value in $[v_m,u_m]$ such that
     \begin{equation}
 \label{q}
 \rho(q_m,Q_1)=u_m,~\rho(q_m,Q_2)=v_m,~\mbox{for all $m\ge1$}.
 \end{equation}
    \item[(iv)] The following inequalities hold:
 \begin{equation}
 \label{diffq1}
 \|q_m-q_n\|\le (\|z\|+2)\left(\max\{u_m,u_n\}-\min\{v_m,v_n\}\right),~\mbox{for all $m,n\ge1$}.
 \end{equation}
\end{enumerate}
 \end{lemma}
 %%%%%%%%%%%%%%%%

%%%%%%%%%%%%%%%%
Now, we are ready to state and prove our main  theorem.  As an improvement of Lemma \ref{lemma1} given  by Borodin in \cite{Borodin1}, this theorem below is the key result in obtaining the equivalence of the (BLT)-property and reflexivity.
\begin{theorem} \label{Borodin}
Let $X$ be an arbitrary infinite-dimensional Banach space, $\{Y_n\}_{n\ge1}$ be a system of strictly nested subspaces with the properties
\begin{enumerate}
\item $\overline Y_n \subset Y_{n+1}$ for all $n\ge1$;
\item For each $n\ge1$, there exists $y_n\in Y_{n+1}\backslash Y_n$ such that
\begin{equation}
\label{Yn}
\rho(y_n,Y_n)=\|y_n\|.
\end{equation}
 \end{enumerate}
 Let  $\{d_n\}_{n\ge1}$ be a sequence of non-negative numbers  decreasing to $0$:
 \begin{equation}
\label{condition'}
 d_n\geq d_{n+1}\ge0, \,\,\mbox{for}~ n \ge 1;~\lim_{n\to\infty}d_n=0 \\
 % \sum_{k = n+1}^\infty d_k~\mbox{for all $n \ge n_0$ at which $d_n > 0$.
 \end{equation}
 and  $$ d_n > \sum_{k = n+1}^\infty d_k~\mbox{for all}\quad n \ge n_0$$
 Then there exists an element $x\in X$ such that $\rho(x,Y_n)=d_n$ for all $n\ge1$ 
\end{theorem}
\begin{proof} We start by assuming that $Y_1 \neq \{0\}$ and  $d_n>0$, for all $n\ge1$. The case $Y_1=\{0\}$ and the case when $d_n=0$ for some $n$ will be discussed separately in the end of this proof.

We define the sequence $\{\tau_j\}_{j\ge1}$ to be
$$
\tau_1=d_1;~\tau_{j}=\min\limits_{k=2,\ldots,j}\{d_{k-1}-d_{k}\},~\mbox{for $j\ge2$}.
$$
In view of the assumption on $\{d_n\}$, we know that the sequence $\{\tau_{j}\}_{j\ge1}$ is non-negative and decreasing to $0$. Since $Y_j\ne \{0\}$, then we can define $Y_0=\{0\}$ and for any integers $j$, $n$ with $1\le j\le n$, we can set
$$
u_n^{(j)}=1+\frac{\tau_n}{2^jd_j},~v_n^{(j)}=1.
$$
\begin{itemize}
\item For $j=1$, in view of Lemma \ref{lemma2}, we can find a sequence $\{q_{1,n}\}_{n\ge 1}\subset Y_{2}\backslash Y_1$ defined by
\begin{eqnarray*}
q_{1,n}&=&v_n^{(1)}(z_1-f_1(z_1)w_1)+\mu_n^{(1)}f_1(z_1)w_1\\
&=&(z_1-f_1(z_1)w_1)+\mu_n^{(1)}f_1(z_1)w_1,
\end{eqnarray*}
such that
\renewcommand{\labelenumi}{\alph{enumi})}
\begin{enumerate}
\item[(i)] $z_1$ is an element in $Y_2\backslash Y_1$ such that $\rho(z_1,Y_0)=2$, $\rho(z_1,Y_1)=1$; $w_1$ is an element in $\overline Y_2$ with $\|z_1-w_1\|=1$.
    \item[(ii)] $f_1:~Y_2\rightarrow \mathbb R$ is a non-zero linear functional such that 
    $$
    f_1(z_1)=\lim\limits_{a\to\infty}\left(a-\frac{\rho(z_1-aw_1,Y_0))}{\rho(w_1,Y_0)}\right).
    $$
    \item[(iii)] $\mu_n^{(1)}$ is some value in $[v_n^{(1)},u_n^{(1)}]=\left[1,1+\frac{\tau_n}{2d_1}\right]$.
    \item[(iv)] We have 
$$
\rho(q_{1,n},Y_{0})=c^{(1)}\left(1+\frac{\tau_n}{2d_1}\right),~\rho(q_{1,n},Y_1)=1,
$$
where $c^{(1)}=1/|1-2f_1(z_1)|$ 
and there exists another constant $c_1>0$ such that
$$
\|q_{1,m}-q_{1,n}\|\le c_1\frac{\tau_m}{2d_1},~\mbox{for all $n\ge m\ge 1$}.
$$
\end{enumerate}
\item For $j=2$, using the condition (\ref{Yn}), we can find $y_2\in Y_3\backslash Y_2$ such that
$$
\rho\left(\frac{y_2}{\|y_2\|},Y_2\right)=1.
$$
Let 
$$
z_2'=\frac{y_2}{\|y_2\|}+z_1,
$$
then $z_2'\in Y_3\backslash Y_2$ and
\begin{equation}
\label{upper:z2}
\|z_2'-z_1\|=\rho(z_2'-z_1,Y_2)=\rho(z_2',Y_2)=1.
\end{equation}
Next we construct $z_2$ as well as $q_{2,n}$. To this end, let
$$
\alpha_{min}=-1/2;~\alpha_{max}=1/2.
$$
By using the triangle inequality and (\ref{upper:z2}) we have
\begin{equation}
\label{alpha_min}
\|z_2'+\alpha_{min}z_1\|\le \|z_2'-z_1\|+|1+\alpha_{min}|\|z_1\|=1+2|1+\alpha_{min}|=2
\end{equation}
and
\begin{equation}
    \label{lower:z2}
    \|z_2'+\alpha_{max} z_1\|\ge|1+\alpha_{max}|\|z_1\|-\|z_2'-z_1\|=2|1+\alpha_{max}|-1=2.
    \end{equation}
Then the continuity of the mapping $\lambda\mapsto \rho(z_2'+\lambda z_1,Y_1)$,  (\ref{alpha_min}), (\ref{lower:z2}) and the intermediate value theorem imply the existence of a real number $\alpha_2\in[-1/2,1/2]$ such that
$$
\rho(z_2'+\alpha_2z_1,Y_0)=2.
$$
Let
$$
z_2=z_2'+\alpha_2z_1,~z_1'=(\alpha_2+1)z_1,
$$ then  $z_2\in Y_3\backslash Y_2$ satisfies $\rho(z_2,Y_0)=2$, $\rho(z_2,Y_2)=1$; and $z_1'$ satisfies $z_1'\in Y_2\backslash Y_0$ and $\|z_2-z_1'\|=1$. We thus can apply Lemma \ref{lemma2} (by taking $(z,w)=(z_2,z_1')$ in the lemma) to obtain, the existence of a sequence $\{q_{2,n}\}_{n\ge 2}\subset Y_{3}\backslash Y_2$ given by:
$$
q_{2,n}=(z_2-f_2(z_2)z_1')+\mu_n^{(2)}f_2(z_2)z_1',
$$
where
\renewcommand{\labelenumi}{\alph{enumi})}
\begin{enumerate}
    \item[(i)] The non-zero linear functional
     $f_2:~Y_3\rightarrow \mathbb R$ is such that 
     $$
     f_2(z_2)=\lim\limits_{a\to\infty}\left(a-\frac{\rho(z_2-az_1',Y_0)}{\rho(z_1',Y_0)}\right)\ge 0.
     $$
    \item[(ii)] $\mu_n^{(2)}$ is some value in $[v_n^{(2)},u_n^{(2)}]=\left[1,1+\frac{\tau_n}{2^2d_2}\right]$ such that
    $$
\rho(q_{2,n},Y_0)=c^{(2)}\left(1+\frac{\tau_n}{2^2d_2}\right),~\rho(q_{2,n},Y_2)=1
$$
for some constant $c^{(2)}>0$.
\item[(iii)] There is a constant $c_2>0$ such that
$$
\|q_{2,m}-q_{2,n}\|\le c_2\frac{\tau_m}{2^2d_2},~\mbox{for all $n\ge m\ge 2$}.
$$
Therefore, by using Lemma \ref{lemma1'},  we get the existence of a non-zero linear functional $f_{1,n}:~Y_{3}\longmapsto\mathbb R$ such that
$$
Y_1\subset \ker f_{1,n},~\|f_{1,n}\|=\frac{1}{\rho(q_{1,n},Y_1)}=f_{1,n}(q_{1,n})=1
$$
and
\begin{eqnarray*}
f_{1,n}(q_{2,n})&\ge&(\mu_n^{(2)}-1)f_2(z_2)-\frac{\rho\left(q_{2,n}-((\mu_n^{(2)}-1)f_2(z_2))q_{1,n},Y_1\right)}{\rho(q_{1,n},Y_1)}.
\end{eqnarray*}
Since
\begin{eqnarray*}
&&\rho\left(q_{2,n}-(\mu_n^{(2)}-1)f_2(z_2)q_{1,n},Y_1\right)
\\
&&\le \rho\left(q_{2,n},Y_1\right)+(\mu_n^{(2)}-1)f_2(z_2)\rho\left(q_{1,n},Y_1\right)\\
&&=\rho\left(q_{2,n},Y_1\right)+(\mu_n^{(2)}-1)f_2(z_2).
\end{eqnarray*}
\end{enumerate}
\end{itemize}
Continuing with  the steps above we can obtain that:
for $j=1,\ldots,n-1$, there exists a sequence $\{q_{j,n}\}_{n\ge j}\subset Y_{j+1}\backslash Y_j$ such that
\begin{equation}
\label{qjnY}
\rho(q_{j,n},Y_{0})=1+\frac{\tau_n}{2^jd_j},~\rho(q_{j,n},Y_j)=1;
\end{equation}
\begin{equation}
\label{diffq}
\|q_{j,m}-q_{j,n}\|\le 4\frac{\tau_m}{2^jd_j},~\mbox{for all $n\ge m\ge j$}.
\end{equation}
Also there is a non-zero linear functional $f_{j,n}:~Y_{j+2}\longmapsto\mathbb R$ such that
$$
Y_j\subset \ker f_{j,n},~\|f_{j,n}\|=f_{j,n}(q_{j,n})=1
$$
and
\begin{equation}
\label{f}
f_{j,n}(q_{j+1,n})\in\left[-1,\frac{\tau_n}{2^{j}d_{j+1}}-1\right].
\end{equation}
Now we fix $n\ge1$. Take $\lambda_{n,n}:=d_n$. We see clearly from (\ref{qjnY}) that
\begin{equation}
\label{identqn}
\rho(\lambda_{n,n}q_{n,n},Y_n)=\lambda_{n,n}\rho(q_{n,n},Y_n)=d_n.
\end{equation}
 Then, first by using (\ref{qjnY}) and (\ref{f}), we obtain
\begin{eqnarray}
\label{upperqn}
&&\rho(\lambda_{n,n}q_{n,n},Y_{n-1})=\lambda_{n,n}\rho(q_{n,n},Y_{n-1})\le \lambda_{n,n}\rho(q_{n,n},Y_0)\nonumber\\
&&=d_n\left(1+\frac{\tau_n}{2^nd_n}\right)\le  d_n+\tau_n\le d_{n-1};
\end{eqnarray}
and, by the properties of $f_{n-1,n}$ in (\ref{f}), we see that
\begin{eqnarray}
\label{lowerqn}
&&\rho(\lambda_{n,n}q_{n,n}-(d_{n-1}+d_nf_{n-1,n}(q_{n,n}))q_{n-1,n},Y_{n-1})\nonumber\\
&&\ge |f_{n-1,n}(\lambda_{n,n}q_{n,n}-(d_{n-1}+d_nf_{n-1,n}(q_{n,n}))q_{n-1,n})|\nonumber\\
&&= -d_nf_{n-1,n}(q_{n,n})+(d_{n-1}+d_nf_{n-1,n}(q_{n,n}))f_{n-1,n}(q_{n-1,n})\nonumber\\
&&=d_{n-1}+d_n\left(f_{n-1,n}(q_{n,n})-f_{n-1,n}(q_{n,n})\right)= d_{n-1}.
\end{eqnarray}
The mapping
$$
\lambda\longmapsto\rho(\lambda_{n,n}q_{n,n}+\lambda q_{n-1,n},Y_{n-1})
$$
 is continuous, through which the image of $[-(d_{n-1}+d_nf_{n-1,n}(q_{n,n})),0]$ contains $d_{n-1}$. Then,  by (\ref{upperqn}), (\ref{lowerqn}) and the intermediate value theorem, there exists $\lambda_{n-1,n}\in [-(d_{n-1}+d_nf_{n-1,n}(q_{n,n})),0]$ such that
\begin{equation}
\label{rhodn}
\rho(\lambda_{n,n}q_{n,n}+\lambda_{n-1,n}q_{n-1,n},Y_{n-1})= d_{n-1}.
\end{equation}
Furthermore, by the fact that $q_{n-1,n}\in Y_n$ and (\ref{identqn}), we have
$$
\rho(\lambda_{n,n}q_{n,n}+\lambda_{n-1,n}q_{n-1,n},Y_{n})= \rho(\lambda_{n,n}q_{n,n},Y_{n})= d_{n}.
$$
For any $2\le k\le n-1$, assume that we can find the real numbers
\begin{eqnarray}
\label{int-lambda}
&&\lambda_{n-1,n}\in[-(d_{n-1}+d_{n}f_{n-1,n}(q_{n,n})),0],\nonumber\\
&&~\ldots,\ldots,\ldots~\nonumber\\
&&\lambda_{k,n}\in[-(d_{k}+d_{k+1}f_{k,n}(q_{k+1,n})),0]
\end{eqnarray}
such that
$$
\rho(\lambda_{n,n}q_{n,n}+\lambda_{n-1,n}q_{n-1,n}+\ldots+\lambda_{k,n}q_{k,n},Y_m)=d_m,~\mbox{for}~m=k,\ldots,n.
$$
Let $
z_{k,n}=\lambda_{n,n}q_{n,n}+\ldots+\lambda_{k,n}q_{k,n}$.
Then, on one hand, by using the triangle inequality, (\ref{int-lambda}), (\ref{qjnY}), (\ref{f}) and  the fact that $k\ge2$, we obtain
\begin{eqnarray}
\label{zY}
&&\rho(z_{k,n},Y_{k-1})\le\sum_{j=k}^n|\lambda_{j,n}|\rho(q_{j,n},Y_{k-1})\nonumber\\
&&\le|\lambda_{n,n}|\rho(q_{n,n},Y_{0})+\sum_{j=k}^{n-1}|\lambda_{j,n}|\rho(q_{j,n},Y_{0})\nonumber\\
&&\le d_n\left(1+\frac{\tau_{n}}{2^nd_n}\right)+\sum_{j=k}^{n-1}\left(d_{j}+d_{j+1}f_{j,n}(q_{j+1,n})\right)\left(1+\frac{\tau_{n}}{2^jd_j}\right)\nonumber\\
&&\le d_n\left(1+\frac{\tau_{n}}{2^nd_n}\right)+\sum_{j=k}^{n-1}\left(d_{j}+d_{j+1}\left(\frac{\tau_n}{2^{j}d_{j+1}}-1\right)\right)\left(1+\frac{\tau_{n}}{2^jd_j}\right)\nonumber\\
&&\le \left(d_n+\sum_{j=k}^{n-1}(d_{j}-d_{j+1})\right)+\sum_{j=2}^{n}\frac{\tau_n(2d_j+d_{j+1}(\tau_n/(2^{j}d_{j+1})-1)}{2^jd_j}\nonumber\\
&&\le d_k+\sum_{j=2}^{n}\frac{\tau_n(2d_j)}{2^{j}d_j}= d_k+\tau_n\sum_{j=2}^n\frac{1}{2^{j-1}}\nonumber\\
&&\le d_k+\tau_n\sum_{j=2}^\infty\frac{1}{2^{j-1}}= d_k+\tau_n\le d_{k-1}.
\end{eqnarray}
On the other hand, by using the properties of $f_{k-1,n}$ in (\ref{f}), we obtain
\begin{equation}
\label{bound:z}
\rho(z_{k,n}-(d_{k-1}+d_kf_{k-1,n}(q_{k-1,n})q_{k-1,n},Y_{k-1})\ge d_{k-1}f_{k-1,n}(q_{k-1,n})= d_{k-1}.
\end{equation}
Therefore, in view of (\ref{zY}), (\ref{bound:z}) and the intermediate value theorem, there is $\lambda_{k-1,n}\in [-(d_{k-1}+f_{k-1,n}(q_{k-1,n})),0]$ such that
$$
\rho(z_{k,n}+\lambda_{k-1,n}q_{k-1,n},Y_{k-1})=d_{k-1}.
$$
Furthermore,
$$
\rho(z_{k,n}+\lambda_{k-1,n}q_{k-1,n},Y_{m})=\rho(z_{k,n},Y_m)=d_{m}~\mbox{for}~m=k,\ldots,n.
$$
Continuing this procedure until $k=1$ is included, we obtain the element
$$
x_{n,n}=\lambda_{n,n}q_{n,n}+\ldots+\lambda_{1,n}q_{1,n},
$$
for which $\rho(x_{n,n},Y_k)=d_k$ and
\begin{equation}
\label{bound_lambdakn}
|\lambda_{k,n}|\left\{
\begin{array}{ll}
\le d_k-d_{k+1}(1-2^{-k})& \mbox{for all $n\ge2$, $k\le n-1$,}\\
=d_n&\mbox{for $n\ge1$, $k=n$.}
\end{array}\right.
 \end{equation}
We see from (\ref{bound_lambdakn}) that for each $k\ge1$, the sequence $\{\lambda_{k,n}\}_{n\ge k}$ is bounded by $d_k$. Then using the usual diagonalization process, we choose a sequence $\Lambda$ of indices $n$ such that, for all $k\ge1$, $\lambda_{k,n}$ converges to the limit $\lambda_k$ as $n\rightarrow\infty$, $n\in\Lambda$. We then claim that $|\lambda_k|\le d_k$ and the limit $\lim\limits_{n\rightarrow\infty}\sum\limits_{k=1}^n\lambda_kq_{k,n}$ exists in $X$. Indeed, by using the triangle inequality, we have for $m\le n$,
\begin{eqnarray}
\label{dfflambdaq}
&&\Big\|\sum_{k=1}^n\lambda_{k}q_{k,n}-\sum_{k=1}^m\lambda_{k}q_{k,m}\Big\|=\Big\|\sum_{k=1}^m\lambda_{k}(q_{k,n}-q_{k,m})+\sum_{k={m+1}}^n\lambda_{k}q_{k,n}\Big\|\nonumber\\
&&\le \sum_{k=1}^m|\lambda_{k}|\|q_{k,n}-q_{k,m}\|+\sum_{k={m+1}}^n|\lambda_{k}|\|q_{k,n}\|.
\end{eqnarray}
First using $|\lambda_k|\le d_k$ and (\ref{diffq}), we obtain
\begin{equation}
\label{dfflambdaq1}
\sum_{k=1}^m|\lambda_{k}|\|q_{k,n}-q_{k,m}\|\le 4\tau_m\sum_{k=1}^m\frac{d_k}{2^kd_k}\xrightarrow[n,m\rightarrow\infty]{}0,
\end{equation}
and  from (\ref{zY}) we see
\begin{equation}
\label{dfflambdaq2}
\sum_{k=m}^n|\lambda_k|\|q_{k,n}\|\le d_{m-1}\xrightarrow[n,m\rightarrow\infty]{}0.
\end{equation}
Combining (\ref{dfflambdaq}), (\ref{dfflambdaq1}) and (\ref{dfflambdaq2}) yields
$$
\left\|\sum_{k=1}^n\lambda_{k}q_{k,n}-\sum_{k=1}^m\lambda_{k}q_{k,m}\right\|\xrightarrow[n,m\rightarrow\infty]{}0,
$$
i.e., $\Big\{\sum\limits_{k=1}^n\lambda_kq_{k,n}\Big\}_{n\ge1}$ is a Cauchy sequence in the Banach space $X$, therefore it has a limit in $X$.
Furthermore, we claim that the element
$$
x:=\lim_{n\rightarrow\infty}\sum_{k=1}^n\lambda_kq_{k,n}
$$
is the limit of the sequence $\{x_{n,n}\}_{n\in\Lambda}$ as $n\to\infty$. By using the fact that
$\|q_{k,n}\|\le 2$, (\ref{bound_lambdakn}) and (\ref{dfflambdaq2}), we obtain
\begin{eqnarray*}
&&\|x-x_{n,n}\|\le \|x-x_n\|+\|x_n-x_{n,n}\|\le \|x-x_n\|+\sum_{k=1}^n|\lambda_{k,n}-\lambda_k|\|q_{k,n}\|\\
&&\le \|x-x_n\|+\sum_{k=1}^{N}|\lambda_{k,n}-\lambda_k|\|q_{k,n}\|+\sum_{k=N+1}^n(|\lambda_{k,n}|+|\lambda_k|)\|q_{k,n}\|\\
&&\le \|x-x_n\|+2\sum_{k=1}^{N}|\lambda_{k,n}-\lambda_k|+d_{n}\|q_{n,n}\|\\
&&+2\sum_{k=N+1}^{n-1}\left(d_k-d_{k+1}+d_{k+1}2^{-k}\right)+\sum_{k=N+1}^n|\lambda_k|\|q_{k,n}\|\\
&&\le \|x-x_n\|+2 N\max_{1\le k\le N,N<n}|\lambda_{k,n}-\lambda_k|+2d_n\\
&&+2\left(d_{N+1}-d_n+d_{N+1}\sum_{k=N+1}^{\infty}2^{-k}\right)+\sum_{k=N+1}^n|\lambda_k|\|q_{k,n}\|\xrightarrow[n\rightarrow\infty,n\in\Lambda,N\rightarrow\infty]{}0.
\end{eqnarray*}
Finally, by the continuity of $x\longmapsto\rho(x,Y_k)$, we obtain
$$
\rho(x,Y_k)=\lim_{n\rightarrow\infty,n\in\Lambda}\rho(x_{n,n},Y_k)=d_k,~\mbox{for all}~k\ge1.
$$
Now we prove the theorem in the following two other particular cases:
\begin{enumerate}
\item If $Y_1=\{0\}$, the problem can be easily converted to the case $Y_1\neq\{0\}$:
\begin{itemize}
\item If $d_1>d_2$, then having constructed an element $z$ with $\rho(z,Y_n)=d_n$ for all $n\ge 2$, we can use Lemma \ref{lemma1} to construct an element $x$ with $\rho(x,Y_k)=d_k$ for $k=1,2$ and such that $x-\lambda z\in Y_{2}$ for some $\lambda>0$. But then observe that
$$
d_{2}=\rho(x,Y_{2})=\rho((x-\lambda z)+\lambda z,Y_2)=\rho(\lambda z,Y_{2})=\lambda d_{2};
$$
therefore, $\lambda=1$ and
$$\rho(x,Y_n)=\rho(z,Y_n)=d_n~\mbox{ for all $n\ge 2$}.
$$
This together with the fact that $\rho(x,Y_1)=d_1$ entails $\rho(x,Y_n)=d_n$ for all $n\ge1$.
\item If $d_1=d_2$, then it suffices to take
$$
q_{1,n}=z_1
$$
with $z_1\in Y_2\backslash Y_1$ satisfying $\|z_1\|=1$. The constructions of
$$q_{2,n},~q_{3,n},\ldots$$ are similar to the above proof.
\end{itemize}
\item If $d_n=0$ starting from some $n$, then the desired element exists by applying the above proof to the setting $X=Y_n$  where the subspaces $Y_1\subset Y_2\subset\ldots\subset Y_{n-1}$  are in $X$.
\end{enumerate}
\end{proof}
%%%%%%%%%%%%%%%%%%%%%%%%%%%%%%%%%
%\section{ Reflexivity and (BLT)-property}
\begin{definition}
\label{def:BLT}
We say a Banach space  $X$ has  the (BLT)-property if it satisfies the above Theorem \ref{Borodin}.
\end{definition}
Recall that a given Banach space  $X$ has the (B)-property if for a sequence of strictly increasing subspaces $\{Y_n\}$ of $X$  with  $\overline{\bigcup_{n=1}^{\infty} Y_n}=X$ and a decreasing null sequence $ \{d_n\}$, there exists an $x\in X$ such that $\rho(x, Y_n) =d_n$ for all $n\in \mathbb{N}$.
 Note that the (BLT)-property is stronger than the (B)-property because of the condition  (\ref{Yn'}) placed on the subspaces. However as we have shown in the introduction this condition is a natural condition.  

Our Theorem \ref{Borodin} yields a stronger equivalence between Banach spaces  which satisfy  the  (BLT)-property  and reflexivity. Indeed, the null sequence $\{d_n\}$ which satisfies the (BLT)-property is not a special sequence like the one satisfying $(B_f)$.

 As we have proved previously  for spaces satisfying the property (B), it is known that if a given Banach space satisfies the (BLT)-property, then it is reflexive. But it was an open problem (see \cite[p. 152 ]{Sin}) for a long time that if we start with a reflexive Banach space $X$  whether or not  $X$ satisfies the (BLT)-property.  Corollary to our main Theorem \ref{Borodin} answers this question in the positive. Let us first recall the following theorem:

\begin{theorem}  \cite[Theorem $2.14$, p. $19$]{Singer2} \label{Sin}
%[See  Theorem $2.14$ on Page $19$ of \cite{Singer2}]\label{Sin}

Let $X$ be a Banach space, then $X$ is reflexive if and only if every closed subspace of  $X$ is proximinal.
\end{theorem}
\begin{corollary}
Let $X$ be a reflexive Banach space then $X$ has the  (BLT)-property.
\end{corollary}
\begin{proof}
Assume  $X$ is reflexive and let $\{Y_n\}$ be a collection of  a closed convex subspaces of $X$ which are strictly nested. Let $z_n \in Y_{n+1} \setminus Y_n$. Then by the above Theorem \ref{Sin},  there exists $v_n \in  Y_n$ such that $d(z_n, Y_n)= \|z_n-v_n\|$. Since $\overline{Y_n} \subset Y_{n+1}$,  we have $y_n= z_n-v_n$ belongs to $Y_{n+1}$ and $d(y_n, Y_n) = \|y_n\|$. Thus conditions of  our main theorem, Theorem \ref{Borodin}  are satisfied and the  (BLT)-property holds true for the reflexive Banach spaces.
\end{proof}
Thus we have the equivalence of reflexive Banach spaces and those Banach spaces which have the (BLT)-property.
\begin{corollary}
\label{corollary:main}
 $X$ is a reflexive Banach space if and only if  $X$ has the  (BLT)-property.
\end{corollary}

\begin{remark}
It is worth noting that any closed convex subset of a reflexive Banach space is proximinal but a closed convex set in general is not proximinal. For example, if we
let $X= \ell_1$, and for any $n \in \mathbb{N}$, $e_n \in X$ be such that its $n$th entry  is $ \frac{n+1}{n}$ and all other entries are $0$, i.e.,
\[\begin{array}{ccccccccc}
e_n & = & (0,&0,&\dots, &\frac{n+1}{n},& 0,&\dots,& 0) \\
&&&&&\uparrow_{nth} &&
\end{array}.\]
 then setting  $K = \overline{\mbox{co}}\{ e_1, e_2, \dots\}$, we observe that $K$ is a closed convex subset of $X$, however $K$ is not proximinal. For more on proximinal sets in Banach spaces we refer the reader to \cite{Ass}.  The classic book of Singer \cite{Sin} contains ample information on proximinal sets as well.
 \end{remark}

%%%%%%%%%%%%%%%%%%%%%%%%%

\begin{remark}%Observe that proximibility implies  the condition (\ref{Yn}) in our main theorem but our condition  (\ref{Yn}) is not equivalent to the  proximinality condition. In the following we make a few observations about proximinal sets.
Let $Y$ be a proper closed subspace of a Banach space $X$, it is known that if $Y$ is proximinal then there exists a norm attaining functional vanishing on $X$. Furthermore,  by the  Bishop-Phelps-Bollobas Theorem \cite{BP} it is known that the set of norm attaining functionals on a Banach space is norm dense in the dual space. It is not difficult to find proximinal hyperplanes since for hyperplanes proximinality and the existence of vanishing norm-attaining functionals  are equivalent.  However, there is  an example given by Charles Read \cite{CR} where he constructs a space in which no subspace of  finite codimension at least $2$ is proximinal.
%%%%%%%%%%%%%%%%%%%%%%%
% It is also well-known that a normed linear space $Y_n$ is reflexive if it is proximinal in  every super space $X$ (see \cite{IS}).  Note also that  if $Y_n$ is proximinal in $Y_{n+1}$ and  $Y_n \subset Y_{n+1} \subset X$ it does not follow directly  that $Y_n$ is proximinal in $X$.  For those special  Banach spaces where this is true, we refer the reader to \cite{EC}.
%For more information on proximinal sets consult the classical book of Singer \cite{Sin}.
\end{remark}
%%%%%%%%
\iffalse
Recall that,  let $X$ be a Banach space and $V$ be its proper closed subspace.
We call $V$ proximinal if for any $x\in X$ there exists $v\in V$ such that
 $$\rho(x, V) =\|x-v\|.$$
 \begin{enumerate}
 \item Hilbert space and finite dimensional Banach space are proximinal.
\item This is also the case when $X$ is a dual space and $V$ is weak* closed.
\item  $M$-ideals are known to be proximinal.
\item Strict convexity does imply proximinality. In fact, any subspace $V$ of a strictly convex $X$ is Chebyshev for any $x\in X$ there exists a unique $v\in V$   such that  $\rho(x, V) =\|x-v\|$.
\end{enumerate}
\fi
%%%%%%%%%%%%
\iffalse
%\end{remark}
\begin{remark}
Necessity of our assumption that $\overline{Y}_n $ is strictly included
in $Y_{n+1}$ can be seen if we consider
%\begin{example}
 $X = L^{\infty}[0,1]$ and $C[0,1] \subset L^{\infty} [0,1]$ and define the subspaces  of $X$ as follows:

\begin{enumerate}
\item
%[$\bullet$]

 $ Y_1 =  \mbox{space of all polynomials}$;
%\item  $Y_2 =$ span$ [Y_1 \cup f_1],$
%where $ f_1 \in C[0,1] \setminus Y_1$,
\item
%[$\bullet$]
 $ Y_{n+1} = $span$[Y_n \cup \{f_{n}\}]$
where $f_{n} \in C[0,1] \setminus Y_n$, for $n\ge1$.
\end{enumerate}
Observe that by the Weierstrass Theorem we have $\overline {Y}_n = C[0,1] $ for all $n \ge1$.
Take any  $ f \in L^{\infty}[0,1]$ and consider the following cases:
\begin{enumerate}
\renewcommand{\labelenumi}{\alph{enumi})}
\item If  $ f \in C[0,1]$ , then
$$
\rho (f,Y_n) = \rho(f,C[0,1]) = 0~\mbox{for all $n\ge1$}.
$$
\item If $ f \in L^{\infty}[0,1] \setminus C[0,1] $, then
$$
\rho(f, Y_n) = \rho(f, C[0,1]) = d > 0~ \mbox{ (independent of $n$)}.
$$
%Note that in the above, we have used the fact that $\rho(f,Y_n) = \rho(f,\overline{Y}_n)$.
Hence in this case BLT does not hold.
\end{enumerate}
\end{remark}
\fi
%%%%%%%%%%%%

\section{Appendix} 
The following two result are used in our theorem.  The first Lemma \ref{lemma1'} is like the classical Hahn-Banach Theorem. Unlike the classical Hahn-Banach Theorem (see e.g. \cite[p. 63]{Dunford}), where the linear functional $f$ is evaluated on one element $x$, in  Lemma \ref{lemma1'} the values of $f$ are evaluated on two elements $x_1,x_2$.
\begin{lemma}
\label{lemma1'}
Let $(X,\|\cdot\|)$ be a normed linear space and let $Q$ be a subspace of $X$ with $\overline Q\subset X$. Pick two elements $x_1,x_2\in X\backslash \overline Q$ such that $x_2\notin \mbox{span} [\{x_1\}\cup Q]$. Then, there exists a non-zero linear functional $f:~X\longmapsto \mathbb R$ such that
\begin{equation}
\label{linearf1}
f(q)=0,~\mbox{for all $q\in Q$},~\|f\|=\frac{1}{\rho(x_1,Q)},~f(x_1)=1
\end{equation}
and
\begin{equation}
\label{linearf2}
f(x_2)=\lim_{a\to+\infty}\left(a-\frac{\rho(x_2-a x_1,Q)}{\rho(x_1,Q)}\right).
\end{equation}
\end{lemma}
\textbf{Proof.} First we note that the existence of $f(x_2)$ in (\ref{linearf2}) is guaranteed by the following 2 facts:
\begin{enumerate}
\item By using the triangle inequality, for all $a\ge0$,
$$
\left|a-\frac{\rho(x_2-a x_1,Q)}{\rho(x_1,Q)}\right|=\frac{\left|a\rho(x_1,Q)-\rho(x_2-a x_1,Q)\right|}{\rho(x_1,Q)}\le\frac{\rho(x_2,Q)}{\rho(x_1,Q)}.
$$
\item Let $a_1\le a_2$, then by using the triangle inequality,
\begin{eqnarray}
\label{increase}
&&\left(a_2-\frac{\rho(x_2-a_2 x_1,Q)}{\rho(x_1,Q)}\right)-\left(a_1-\frac{\rho(x_2-a_1 x_1,Q)}{\rho(x_1,Q)}\right)\nonumber\\
&&=(a_2-a_1)-\frac{\rho(x_2-a_2 x_1,Q)-\rho(x_2-a_1 x_1,Q)}{\rho(x_1,Q)}\nonumber\\
&&\ge (a_2-a_1)-\frac{\rho((a_2-a_1) x_1,Q)}{\rho(x_1,Q)}\nonumber\\
&&=0.
\end{eqnarray}
\end{enumerate} 
Now,  let $U$ be the linear subspace of $X$ spanned by $\{x_1\}\cup Q$, then any element in $U$ has unique decomposition of the form $q+\alpha x_1$, for some $q\in Q$, $\alpha\in\mathbb R$. From Hahn-Banach Theorem, we have  the linear functional $g:~U\longmapsto \mathbb R$ defined by
\begin{equation}
\label{def:g}
g(q+\alpha x_1)=\alpha,~\mbox{for any $q\in Q$, $\alpha\in\mathbb R$}
\end{equation}
satisfies
\begin{equation}
\label{linearg}
g(q)=0,~\mbox{for all $q\in Q$},~\|g\|=\frac{1}{\rho(x_1,Q)},~\mbox{and}~g(x_1)=1.
\end{equation}
Since $x_2\notin U$, we then let $S$ be the linear subspace of $X$, spanned by $\{x_2\}\cup U$. Define the \emph{sublinear functional} $p$ to be
\begin{equation}
\label{px}
p(x)=\|g\|\|x\|,~\mbox{for $x\in S$}.
\end{equation}
We then can extend $g$ to a linear functional $\tilde g:~S\longmapsto \mathbb R$, by taking
\begin{equation}
\label{def:nu}
\tilde g(q+\alpha x_1+\beta x_2)=\alpha+\beta \nu,~\mbox{for all $q\in Q$ and $\alpha,\beta\in\mathbb R$},
\end{equation}
where the real number $\nu$ is chosen so that
\begin{equation}
\label{fp}
\tilde g(x)\le p(x),~\mbox{for $x\in S$}.
\end{equation}
From Hahn-Banach Theorem (see \cite{Dunford}, Page 63, inequality (i)), any $\nu$ satisfying the inequalities (\ref{inequality1}) below yields (\ref{fp}).
\begin{equation}
\label{inequality1}
\left\{\begin{array}{ll}
&\nu\ge\max\limits_{q\in Q,~\alpha\in \mathbb R}\left\{-p(-q-\alpha x_1-x_2)-g(q+\alpha x_1)\right\};\\
&\nu\le \min\limits_{q'\in Q,~\alpha'\in \mathbb R}\left\{p(q'+\alpha' x_1+x_2)-g(q'+\alpha' x_1)\right\}.
\end{array}\right.
\end{equation}
Next, we show that $\nu$ can take the value $\lim\limits_{a\to+\infty}\left(a-\frac{\rho(x_2-a x_1,Q)}{\rho(x_1,Q)}\right)$. Using (\ref{def:g}), we know that it is enough to prove 
\begin{equation}
\label{inequality2}
-p(-q-\alpha x_1-x_2)-\alpha\le \lim\limits_{a\to+\infty}\left(a-\frac{\rho(x_2-a x_1,Q)}{\rho(x_1,Q)}\right)\le p(q'+\alpha' x_1+x_2)-\alpha',
\end{equation}
for all $q,q'\in Q$ and $\alpha,\alpha'\in\mathbb R$. Equivalently, we will show
\begin{equation}
\label{inequality21}
 \lim\limits_{a\to+\infty}\left(a-\frac{\rho(x_2-a x_1,Q)}{\rho(x_1,Q)}\right)\ge-p(-q-\alpha x_1-x_2)-\alpha,~\mbox{for $q\in Q$ and $\alpha\in\mathbb R$}
\end{equation}
and
\begin{equation}
\label{inequality22}
\lim\limits_{a\to+\infty}\left(a-\frac{\rho(x_2-a x_1,Q)}{\rho(x_1,Q)}\right)\le p(q'+\alpha' x_1+x_2)-\alpha',~\mbox{for $q'\in Q$ and $\alpha'\in\mathbb R$}.
\end{equation}
% \begin{description}
% \item[(1)]
%\begin{enumerate}
 To show (\ref{inequality21}), we combine (\ref{px}), (\ref{linearg}),  Property 3 of  $\rho( \cdot , Q)$, the fact that $q\in Q$ and (\ref{increase}) to obtain:
      \begin{eqnarray}
     \label{uppermu}
     &&-p(-q-\alpha x_1-x_2)-\alpha=-\frac{\|q+\alpha x_1+x_2\|}{\rho(x_1,Q)}-\alpha\nonumber\\
     &&\le -\frac{\rho(q+\alpha x_1+x_2,Q)}{\rho(x_1,Q)}-\alpha=-\frac{\rho(\alpha x_1+x_2,Q)}{\rho(x_1,Q)}-\alpha\nonumber\\
     &&\le \lim_{a\to+\infty}\left(a-\frac{\rho(x_2-a x_1,Q)}{\rho(x_1,Q)}\right)
     \end{eqnarray}
     is true  for all $q\in Q$ and $\alpha\in \mathbb R$.

      To show (\ref{inequality22}),  again we apply (\ref{px}) and Property 3  of $\rho(\cdot , Q)$, to obtain for all $q'\in Q$ and $\alpha'\in \mathbb R$,
     \begin{equation}
     \label{lowermu}
     p(q'+\alpha' x_1+x_2)-\alpha'=\frac{\|q'+\alpha' x_1+x_2\|}{\rho(x_1,Q)}-\alpha'\ge\frac{\rho(\alpha' x_1+x_2,Q)}{\rho(x_1,Q)}-\alpha'
     \end{equation}
     For each $\alpha'\in\mathbb R$, there is  $a>0$ large enough such that
     \begin{eqnarray}
\label{increase_1}
&&\left(a-\frac{\rho(x_2-a x_1,Q)}{\rho(x_1,Q)}\right)-\left(-\alpha'+\frac{\rho(x_2+\alpha' x_1,Q)}{\rho(x_1,Q)}\right)\nonumber\\
&&=(a+\alpha')-\frac{\rho(x_2-a x_1,Q)+\rho(x_2+\alpha' x_1,Q)}{\rho(x_1,Q)}\nonumber\\
&&\le (a+\alpha')-\frac{\rho((a+\alpha') x_1,Q)}{\rho(x_1,Q)}\nonumber\\
&&=0.
\end{eqnarray}

     Therefore combining (\ref{lowermu}) and (\ref{increase_1}), we obtain  (\ref{inequality22}).

 Having proved (\ref{inequality21}) and (\ref{inequality22}), we see  that (\ref{inequality2}) holds. Thus we can take $$
 \nu=\lim_{a\to+\infty}\left(a-\frac{\rho(x_2-a x_1,Q)}{\rho(x_1,Q)}\right)
 $$ in (\ref{def:nu}), to obtain (\ref{linearf2}).

 Now, we show $\|\tilde g\|=\|g\|=\displaystyle \frac{1}{\rho(x_1,Q)}$. This is true, because  the fact that $\tilde g(q)=g(q)$ for $q\in Q$ implies $\|\tilde g\|\ge\|g\|$; and the fact that $\tilde g(x)\le p(x)=\|g\|\|x\|$ leads to $\|\tilde g\|\le \|g\|$. It then follows from (\ref{linearg}) that $\|\tilde g\|=\|g\|=\displaystyle \frac{1}{\rho(x_1,Q)}$.

 Finally, from the Hahn-Banch Extension Theorem $\tilde g$ can be further extended to a linear functional $f:~X\longmapsto \mathbb R$, that satisfies $f(x)=\tilde g(x)$ for $x\in S$ and $\|f\|=\|\tilde g\|$.  $\square$
%%%%%%%%%%%%%%%%%

  The following Lemma \ref{lemma2} is another technical statement. It reveals that, small changes of  $\{d_n\}$ in the Bernstein Lethargy Theorem,  can lead to small position changes of the corresponding element $x\in X$.
 \begin{lemma}
 \label{lemma2}
 Let $Q_1,Q_2$ be two subspaces of an arbitrary normed linear space $(Q_3,\|\cdot\|)$, such that $\overline Q_k\subset Q_{k+1}$ for $k=1,2$. Also assume that for $k=1,2$ there is $y_k\in Q_{k+1}\backslash  Q_k$ for which $\rho(y_k,Q_k)=\|y_k\|$. Let $\{u_m\}_{m\ge1}$ and $\{v_m\}_{m\ge1}$ be two sequences of non-negative numbers, with $u_m\ge v_m$ for all $m\ge1$. Then there exists a sequence of elements  $\{q_m\}_{m\ge1}\subset Q_{3}$, given by
$$
q_{m}=v_m(z-f(z)w)+\mu_mf(z)w,~\mbox{for $m\ge1$},
$$
such that
\begin{enumerate}
\renewcommand{\labelenumi}{\alph{enumi})}
\item [(i)]$z$ is an element in $Q_3\backslash \overline Q_2$ satisfying $\rho(z,Q_1)=\rho(z,Q_2)=1$; $w$ is an element in $\overline Q_3\backslash\{0\}$ such that $\|z-w\|=1$.
    \item[(ii)] $f:~Q_3\rightarrow \mathbb R$ is a non-zero linear functional such that 
    $$
    f(z)=\lim_{a\to+\infty}\left(a-\frac{\rho(z-aw,Q_1)}{\rho(w,Q_1)}\right).
    $$
    \item[(iii)] For each integer $m\ge1$, the variable $\mu_m$ in the expression of $q_m$ is some value in $[v_m,u_m]$ such that
     \begin{equation}
 \label{q}
 \rho(q_m,Q_1)=u_m,~\rho(q_m,Q_2)=v_m,~\mbox{for all $m\ge1$}.
 \end{equation}
    \item[(iv)] The following inequalities hold:
 \begin{equation}
 \label{diffq1}
 \|q_m-q_n\|\le (\|z\|+2)\left(\max\{u_m,u_n\}-\min\{v_m,v_n\}\right),~\mbox{for all $m,n\ge1$}.
 \end{equation}
\end{enumerate}
 \end{lemma}
 %%%%%%%%%%%%%%%%%%%%
\begin{proof} We show (i)-(iv) hold one by one.
  \begin{description}
  \item[\emph{Proof of (i)}] Construction of $z,w$.

By assumption there is $y_1\in Q_3\backslash  Q_2$ such that $\rho(y_1,Q_2)=\|y_1\|$. Taking $y_2=y_1/\|y_1\|$, we obtain $\rho(y_2,Q_2)=\|y_2\|=1$. This together with the fact that $\overline Q_1\subset Q_2$ yields
$$
1=\rho(y_2,Q_2)\le\rho(y_2,Q_1)\le1.
$$
Hence
$$
\rho(y_2,Q_1)=1.
$$
For $x\in\overline Q_2\backslash Q_1$, we set
\begin{equation}
    \label{def:z}
    z=y_2+x\in Q_3~\mbox{and}~w=2z- x\in Q_3.
\end{equation}
Then we have
 \begin{eqnarray}
 \label{zQ}
 &&\rho(z,Q_2)=\rho(y_2+ x,Q_2)=\rho(y_2,Q_2)=1;\nonumber\\
 &&\rho(z,Q_1)=1;\nonumber\\
 &&\|z- x\|=\|y_2\|=1.
 \end{eqnarray}
 By definition of $w$ in (\ref{def:z}) we also get
 \begin{equation}
 \label{z-wQ1}
 \rho(w,Q_1)\ge \rho(2z-x,Q_2)= \rho(2z,Q_2)=2
 \end{equation}
 and
 \begin{equation}
 \label{z-wQ2}
 \rho(z-w,Q_1)\le \|z-w\|=\|z-x\|=1.
 \end{equation}
Hence $\emph{\bf{(i)}}$ is proved.
\item[\emph{Proof of (ii)}] Construction of $f$.

Since $z=(w+x)/2$ with $x\in \overline Q_2\backslash Q_1$, $z\notin \mathcal S= \mbox{span}[\{w\}\cup Q_1]$, then we can apply Lemma \ref{lemma1'} to confirm the existence of  a non-zero real-valued linear functional $f:~Q_3\rightarrow \mathbb R$ such that
\begin{eqnarray}
\label{f1}
&&Q_1\subset \ker f,~\|f\|=\frac{1}{\rho(w,Q_1)},~f(w)=1,\nonumber\\
&&f(z)=\lim_{a\to+\infty}\left(a-\frac{\rho(z-a w,Q_1)}{\rho(w,Q_1)}\right).
\end{eqnarray}
(ii) is proved.
\item[\emph{Proof of (iii) and (iv)}] Construction of the sequence $\{q_m\}_{m\ge1}$.


Let $f$ be the linear functional given in (\ref{f1}). We define
 \begin{equation}
 \label{x}
x_1=f(z)w,~x_2=z-f(z)w.
 \end{equation}
 We will show that the sequence $\{q_m\}_{m\ge1}$ satisfying (\ref{q}) and (\ref{diffq1}) can be found in $\mbox{span}[\{x_1,x_2\}]$. Using (\ref{x}) and (\ref{zQ}) we obtain
\begin{eqnarray*}
 &&\rho(x_2,Q_2)= \rho(z-f(z)w,Q_2)\\
 &&=\rho (z-f(z)(2z-x), Q_2)=\left|1-2f(z)\right|\rho(z,Q_2)\\
 &&=|1-2f(z)|.
 \end{eqnarray*}
 Therefore
 \begin{equation}
 \label{upper1}
 \rho(v_mx_2,Q_2)=v_m\left|1-2f(z)\right|.
 \end{equation}
 On one hand, by using (\ref{x}), (\ref{z-wQ1})  and the fact that $v_m\le u_m$, we have
 \begin{equation}
 \label{upperu}
\rho(v_mx_2+v_mx_1,Q_1)=|v_m|\rho(z,Q_1)=v_m\le u_m.
 \end{equation}
 On the other hand, by using the fact that the kernel $\ker f$ satisfies
 \begin{equation}
 \label{kerf}
 \rho(x,\ker f)=\frac{|f(x)|}{\|f\|},~\mbox{for all $x\in Q_3$},
 \end{equation}
 and the equations (\ref{x}) and (\ref{f1}), we obtain
\begin{eqnarray}
\label{loweru'}
&&\rho(v_mx_2+u_mx_1,Q_1)\ge \rho(v_mx_2+u_mx_1,\ker f)=\frac{|f(v_mx_2+u_mx_1)|}{\|f\|}\nonumber\\
&&=\rho(w,Q_1)\left|v_mf(z-f(z)w)+u_mf(f(z) w)\right|\nonumber\\
&&=\rho(w,Q_1)u_m\left|f(z)\right|.
\end{eqnarray}
From (\ref{f1}), the facts that $\rho(z-w,Q_1)\le1$ (see (\ref{z-wQ2})), we observe that
\begin{equation}
\label{f_control}
\rho(w,Q_1)f(z)\ge \rho(w,Q_1)\left(1-\frac{\rho(z-w,Q_1)}{\rho(w,Q_1)}\right)\ge\rho(w,Q_1)-1.
\end{equation}
It then follows from (\ref{loweru'}), (\ref{f_control}) and the fact that $\rho(w,Q_1)\ge2$ that
\begin{equation}
\label{loweru}
\rho(v_mx_2+u_mx_1,Q_1)\ge \left(\rho(w,Q_1)-1\right)u_m\ge  u_m.
\end{equation}
Since the mapping $\lambda \longmapsto \rho(v_mx_2+\lambda x_1,Q_1)$ is continuous, it follows from (\ref{upperu}), (\ref{loweru})  and the intermediate value theorem that there is a real number $\mu_m\in[v_m,u_m]$ such that
$$
\rho(v_mx_2+\mu_m x_1,Q_1)=u_m~\mbox{and}~\rho(v_mx_2+\mu_m x_1,Q_2)=cv_m.
$$
We then denote $\{q_m\}$ to be
\begin{equation}
\label{q1}
q_m=v_mx_2+\mu_m x_1,~\mbox{for all $m\ge1$}.
\end{equation}
As a consequence (\ref{q}) holds.

Now we show (\ref{diffq1}) holds.  To this end we first state the following two evident facts:
\begin{enumerate}
\item If $\mu_m\in[v_m,u_m]$ for any $m\ge1$, then we have
\begin{equation}
\label{diffmu}
|\mu_m-\mu_n|\le \max\{u_m,u_n\}-\min\{v_m,v_n\},~\mbox{for any $m,n\ge1$}.
\end{equation}
\item For any four real numbers $u_m,v_m,u_n,v_n$ such that $u_m\ge v_m,~u_n\ge v_n$, the following inequality holds:
\begin{equation}
\label{vmvn}
|v_m-v_n|\le \max\{u_m,u_n\}-\min\{v_m,v_n\}.
\end{equation}
\end{enumerate}
It results from (\ref{q1}), the triangle inequality, (\ref{diffmu}) and (\ref{vmvn}) that for any $m,n\ge1$,
\begin{eqnarray*}
\label{dq}
\|q_m-q_n\|&\le& |v_m-v_n|\|z-f(z)w\|+|\mu_m-\mu_n||f(z)|\|w\|\nonumber\\
&\le& c\left(\max\{\mu_m,\mu_n\}-\min\{v_m,v_n\}\right)\nonumber\\
&\le& c\left(\max\{u_m,u_n\}-\min\{v_m,v_n\}\right),
\end{eqnarray*}
where
$$
c=\max\left\{\|z-f(z)w\|,|f(z)|\|w\|\right\}.
$$
Below we show that  $c$ has an upper bound: 
\begin{eqnarray*}
c&=& \max\left\{\|z-f(z)w\|,|f(z)|\|w\|\right\}\\
&\le& \|z\|+|f(z)|\|w\|\\
&\le& \|z\|+\|f\|\|z\|\|w\|\\
&\le& \|z\|+\frac{\|z\|\|w\|}{2}.
\end{eqnarray*}

Hence (\ref{diffq1}) holds. (iii), (iv) are then proved by combining (\ref{q}) and (\ref{diffq1}). The proof of Lemma \ref{lemma2} is completed.
\end{description}
\end{proof}
\fi
%%%%%%%%%%%%%%%%

%%%%%%%%%%%%%%%
\noindent \textbf{Acknowledgement.} Authors would like to thank Jose Mar\'ia Almira for his valuable comments and to Yu. A. Skvortsov for his help in locating the reference \cite{Nik}.

%\bibliographystyle{amsplain}
\begin{thebibliography}{00}

%\bibitem{haag1} U. Haagerup, \textit{Solution of the similarity problem for cylic representations
%of $C^*$-algebras}, Ann. of Math. (2) {\bf 118} (1983), no. 2,
%215--240.
\bibitem{Aksoy} A.G. Aksoy, M. Al-Ansari, C. Case, Q. Peng, Subspace condition for Bernstein's lethargy theorem, Turk J. Math. 41(5) (2017) 1101--1107.
\bibitem{Ak-Al} A.G. Aksoy,  J.M. Almira, On Shapiro's lethargy theorem and some applications, Ja\'{e}n J. Approx. 6(1) (2014) 87--116.
    \bibitem{Ak-Le2}  A.G. Aksoy, G. Lewicki, Bernstein's lethargy theorem in Fr\'{e}chet spaces,  J. Approx. Theory 209 (2016) 58--77.
        \bibitem{Ak-Le} A.G. Aksoy, G. Lewicki,  Diagonal operators, $s$-numbers and Bernstein pairs,  Note Mat. 17 (1999) 209--216.
        \bibitem{AL} {J. M. Almira, U. Luther, } Compactness and generalized approximation spaces. Numerical Functional Analysis and Optimization, 23 (1--2), (2002) 1--38. 
     %   \bibitem{Ak-Pe2} A.G. Aksoy, Q. Peng, Constructing an element of a Banach space with given deviation from its nested subspaces, Khayyam J. Math. 4(1) (2018) 59--76.
\bibitem{Alb} G. Albinus, Remarks on a theorem of S.N. Bernstein, Studia Math. 38 (1970) 227--234.
\bibitem{Al-To} J.M. Almira, N. del Toro, Some remarks on negative results in approximation theory, in:  Proceedings of the Fourth International Conference on Functional Analysis and Approximation Theory, Vol. I (Potenza, 2000). Rend. Circ. Mat. Palermo (2) Suppl. 2002, no. 68, part I, 245--256.
%\bibitem{JMA}  J.M. Almira, U. Luther, Compactness and generalized approximation spaces, Numer. Funct. Anal. Optimiz., 23(1\&2) (2002) 1--38.

\bibitem{Al-Oi} J.M. Almira, T. Oikhberg, Approximation schemes satisfying Shapiro's theorem, J. Approx. Theory 164(5) (2012) 534--571.

%\bibitem{Ass} A. Assadi, H. Haghshenas,  T.D. Narang, A look at proximinal and Chebyshev sets in Banach spaces, Le Matematiche, Vol. LXIX (Fasc. I) (2014) 71--78.
\bibitem{Bernstein} S.N. Bernstein, On the inverse problem in the theory of best approximation of continuous functions, Collected works (in Russian), vol II, Izd.  Akad. Nauk, USSR, 1954, pp 292--294.
%\bibitem{BP}  E. Bishop, R.R. Phelps, A proof that every Banach space is supreflexive, Bull. Amer. Math. Soc. 67 (1961) 97--98.
\bibitem{Borodin1} P.A. Borodin, On the existence of an element with given deviations from an expanding system of
subspaces, Math. Notes 80 (5) (2006) 621--630 (translated from Matematicheskie Zametki
80 (5) (2006) 657--667).
%\bibitem{EC} E. W. Cheney, \textit{The embeddings of proximinal sets}, J. Approx. Theory 48, (1986), no2, 213-225.

\bibitem{De-Hu} F. Deutsch, H. Hundal, A generalization of Tyuriemskih's lethargy theorem  and some applications, Numer. Func. Anal. Opt. 34(9) (2013) 1033--1040.

% \bibitem{Dunford} N. Dunford, J. Schwartz, Linear Operators, Part 1, General Theory, Wiley Classics Library, 1988.
 \bibitem{Hir} R.A. Hirschfeld, On best approximations in normed vector spaces. Nieuw Arch. Wisk (3) 6 (1958), 41-51.
% \bibitem{Hut}  C. Hutton, J.S.  Morrell, J.R. Retherford,  Diagonal operators, approximation numbers and Kolmogoroff diameters, J. Approx. Theory  16  (1976) 48--80.
\bibitem{RCJ} R.C. James, Reflexivity and the supremum of linear functionals,  Ann. of Math. 66(1) (1957) 159--169.
%\bibitem{Konyagin} S.V. Konyagin,  Deviation of elements of a Banach space from a system of subspaces, in: Proceedings of the Steklov Institute of Mathematics 284, no. 1, 2014, pp 204--207.
%    \bibitem{Lew} G. Lewicki, Bernstein's ``lethargy" theorem in metrizable topological linear spaces, Mh. Math. 113 (1992) 213--226.
    \bibitem{Lew1} G. Lewicki, A theorem of Bernstein's type for linear projections, Univ. Jagiellonian. Acta Math. 27 (1988), 23--27.
\bibitem{Mich} B. Micherda, Bernstein's lethargy theorem in SF-spaces, Z. Anal. Anwendungen 22(1) (2003) 3--16.
\bibitem{Nik} V.N. Nikol'skii, Some properties of reflexive spaces (Russian), Kalinin. Gos. Ped. Inst. Ucen. Zap. 29 (1963) 121--125.
\bibitem{Ple1} W. Ple\'{s}niak, On a theorem of S.N. Bernstein in $F$-Spaces, Zeszyty Naukowe Uniwersytetu Jagiellonskiego, Prace Mat. 20 (1979) 7--16.
\bibitem{Ple2} W. Ple\'{s}niak, Quasianalytic functions in the sense of Bernstein, Dissertations Math. 147 (1977) 1--70.
%\bibitem {CR} C.J. Read, Banach spaces with no proximinal subspaces of codimension 2,  ISR. J. Math. (2017). https://doi.org/10.1007/s11856-017-1627-3.
\bibitem{Sha} H.S. Shapiro, Some negative theorems of approximation theory, Michigan Math. J. 11 (1964) 211--217.
\bibitem{Singer2} I. Singer, The Theory of Best Approximation and Functional Analysis, Regional Conference Series in Applied Mathematics, Vol. XIII,  SIAM, Philadelphia, 1974.
%\bibitem {IS} I. Singer,\textit{On normed linear spaces which are proximinal in every super space}, J. Approx. Theory 7, (1973), 399-402.
\bibitem{Sin} I. Singer, Best Approximation in Normed Linear Spaces by Elements of Linear Subspaces, Springer-Verlag, Berlin, 1970.
\bibitem{Tim} A.F. Timan, Theory of Approximation of Functions of a Real Variable, Dover, New York, 1994.
\bibitem{Tyu} I.S. Tyuriemskih, On one problem of S.N. Bernstein, in: Scientific Proceedings of Kaliningrad State Pedagog. Inst. 52,  1967, pp 123--129.
\bibitem {Tyu2}  I.S. Tyuriemskih, The $B$-property of Hilbert spaces,  Uch. Zap. Kalinin. Gos. Pedagog. Inst. 39 (1964)  53--64.
\bibitem{Vas} A.I. Vasil'ev, The inverse problem in the theory of best approximation in $F$-spaces, (In Russian) Dokl. Ross. Akad. Nauk. 365(5) (1999),  583--585.
%\bibitem{MUR} G. J. Murphy, \textit{$C^*$-Algebras and Operator Theory}, Academic Press, Boston, 1990.

%\bibitem{MM} M. Mirzavaziri and M. S. Moslehian, \textit{Automatic continuity of $\sigma$-derivations in
%$C^*$-algebras}, Proc. Amer. Math. Soc. \textbf{134} (2006), no. 11,
%3319--3327.

%\bibitem{H} M. S. Moslehian, \textit{Ky Fan inequalities}, Linear Multilinear Algebra, arXiv:1108.1467v2 (to appear).

%\bibitem{RAS} Th. M. Rassias, \textit{Stability of the generalized orthogonality functional equation},
%Inner product spaces and applications, 219--240, Pitman Res. Notes
%Math. Ser., 376, Longman, Harlow, 1997.

%\bibitem{BRV} J. Bichon, A. De Rijdt, and S. Vaes, \textit{Ergodic coactions with large multiplicity and monoidal equivalence of quantum groups}, Comm. Math. Phys. \textbf{262} (2006), no. 3, 703--728.

\end{thebibliography}

\end{document}
%------------------------------------------------------------------------------
% End of journal.tex
%------------------------------------------------------------------------------




